\subsection{The cokernel of the reduced norm}\label{sec:second-three-coker}
In this subsection, we compute the image and the cokernel of the reduced norm on $K_2$ over three-dimensional local fields (Theorem~\ref{thm:second-three-coker}).
For a three-dimensional local field $F$, we write $\inv_F\colon H^4(F,\Q/\Z(3))\xrightarrow{\sim}\Q/\Z$ for Kato's isomorphism \cite[Theorem 2]{KatoExistence}, as in the Conventions.
For a three-dimensional local field $F$ and a central division algebra $D$ over $F$ with $\alpha=[D]$, duality reduces the computation to showing that a Brauer class that vanishes over every finite splitting field of $\alpha$ lies in the subgroup generated by $\alpha$.

For a field $F$ and $z\in\Br(F)$, let $B_z(F)$ denote the common kernel in $\Br(F)$ of restriction to all finite splitting fields of $z$.
For the primary decomposition $z=\sum_\ell z_\ell$, one has $B_z(F)=\bigoplus_\ell B_{z_\ell}(F)$.
Indeed, $B_z(F)$ is killed by $\ind(z)$.
Over a splitting field of $z_\ell$, we split the other components in degree prime to $\ell$ by successively adjoining separable maximal subfields for them; restriction and corestriction then show that the $\ell$-component lies in $B_{z_\ell}(F)$.
The reverse inclusion follows from the definition.
The \emph{Witt part} is the subgroup of Brauer classes killed by the maximal unramified extension with respect to the top discrete valuation of $F$, whose residue field is the first residue field $k$.
A class in the Witt part has a Witt decomposition $i(a)+(\chi,\pi)$ with $a\in\Br(k)$ and $\chi\in H^1(k,\Q/\Z)$; we call it purely unramified if $\chi=0$.
In first residue characteristic $p$, we call a $p$-primary class wild if it does not belong to the Witt part.
For a two-dimensional local field $F$, the equality $B_z(F)=\langle z\rangle$ follows from Kato's degree-one norm theorem and duality, since the sum of $K_1$ norms from all splitting fields is the reduced-norm image of the division algebra in the Brauer class $z$ \cite[\S3.1, Theorem 3(2)]{KatoLocalII}, \cite[Theorem 2]{KatoExistence}.
Here, for the division algebra $D$ in the class $z$, the group $K_1(D)$ is generated by transfers from separable splitting fields: elements generating separable maximal subfields form a nonempty open subset of $D^\times$, and every element is a ratio of two of them.


\begin{lemma}\label{lem:second-three-norm-quotient}
Let $F$ be a local field of dimension $h\ge2$, let $\ell$ be a prime, and let $L/F$ be a cyclic extension of degree $\ell$.
Then $F^\times/N_{L/F}L^\times$ is an $\F_\ell$-vector space of dimension at least two.
\end{lemma}
\begin{proof}
We use $v_F(N_{L/F}x)=f(L/F)\,v_L(x)$.
The quotient is killed by $\ell$.
Let $k$ be the first residue field.
If the extension is unramified, there are the valuation quotient $\Z/\ell$ and the residue norm quotient of units.
For the cyclic residue extension $l/k$, if the residue norm quotient of units vanished, the projection formula would make the Milnor norm $K_{h-1}^M(l)\to K_{h-1}^M(k)$ surjective, whereas its cokernel is $\Gal(l/k)\simeq\Z/\ell$ by higher local class field theory \cite[\S0, p.\ 168]{KatoExistence}.
If the extension is totally ramified, residues of unit norms lie in $k^{\times\ell}$.
A preimage of a unit norm is a unit, so the images of units in the norm quotient surject onto $k^\times/k^{\times\ell}$.
When $\ell\ne\operatorname{char}k$, tame cyclic total ramification implies $\mu_\ell\subset k$.
Reduction preserves the orders of roots of unity at invertible primes, and at a mixed-characteristic step these orders are bounded by $v(\zeta_{\ell^r}-1)=v(\ell)/(\ell^{r-1}(\ell-1))$.
Thus $k$ has a root of unity of maximal $\ell$-power order.
Such a root and a uniformizer give independent classes in $k^\times/k^{\times\ell}$.
If $\ell=\operatorname{char}k$, the classes of a uniformizer $w$ and $1+w$ are independent.
Indeed, valuation first removes the exponent of $w$.
For a nonzero power of $1+w$ with exponent less than $\ell$, subtracting one gives valuation one, but subtracting one from an $\ell$-th power gives valuation divisible by $\ell$.
The remaining case has ramification index one and purely inseparable residue degree $\ell$ (here $\ell=\operatorname{char}k$).
If a norm is a first principal unit, the residue of its preimage is one, so the preimage is also a first principal unit.
The norm of the preimage lies in the second principal units: for a common uniformizer $\pi$ one has $N_{L/F}(1+\pi a)\equiv1+\pi\Tr_{l/k}(\bar a)\pmod{\pi^2}$, and $\Tr_{l/k}=0$ for the purely inseparable residue extension $l/k$.
Thus $N_{L/F}L^\times\cap U_F^1\subseteq U_F^2$, so the image of the principal units in the norm quotient surjects onto $U_F^1/U_F^2\simeq(k,+)$.
Since $k$ is infinite, the dimension is at least two.
These cases exhaust the possibilities because a complete discretely valued field satisfies $[L:F]=e(L/F)f(L/F)$.
\end{proof}

\begin{lemma}\label{lem:second-three-cyclic-family}
Let $\ell$ be a prime, let $F$ be a local field of dimension at least two, and let $\eta\in\Br(F)\{\ell\}$ and $\theta\in H^1(F,\Q_\ell/\Z_\ell)$.
If $\theta$ has order $n$, then
\[
 \bigcap_{u\in F^\times}\langle\eta+(\theta,u)\rangle
   =\langle n\eta\rangle .
\]
\end{lemma}
\begin{proof}
The case $n=1$ is clear, so assume $n>1$.
By Lemma~\ref{lem:second-three-norm-quotient} and the description of cyclic relative Brauer groups applied to the character $(n/\ell)\theta$ of order $\ell$, the group $(n/\ell)(\theta,F^\times)$ also has dimension at least two.
We may therefore choose $\delta=(\theta,u)$ with $(n/\ell)\delta\notin\langle\eta\rangle$.
Then $\delta$ has order $n$ and $\langle\delta\rangle\cap\langle\eta\rangle=0$, since every nonzero subgroup of the cyclic $\ell$-group $\langle\delta\rangle$ contains $(n/\ell)\delta$; hence
\[
 \langle\eta\rangle\cap\langle\eta+\delta\rangle
     =\langle n\eta\rangle .
\]
The terms for $u=1$ and this $u$ give one inclusion, and the other follows because every $(\theta,u)$ is killed by $n$.
\end{proof}

\begin{lemma}\label{lem:second-three-witt-common-kernel}
Let $F$ be a complete discretely valued field whose first residue field $k$ is an $h$-dimensional local field with $h\ge2$.
For a prime $\ell$, assume $B_z(k)=\langle z\rangle$ for every $z\in\Br(k)\{\ell\}$.
If $\alpha\in\Br(F)\{\ell\}$ is split by a finite unramified extension, then $B_\alpha(F)=\langle\alpha\rangle$.
\end{lemma}
\begin{proof}
Restriction and corestriction through a separable maximal subfield make $B_\alpha(F)$ $\ell$-primary.
Every $\beta$ in this group is killed by the same finite unramified splitting field as $\alpha$, so write the Witt decompositions
\[
 \alpha=i(a)+(\chi,\pi),\qquad\beta=i(b)+(\psi,\pi).
\]
Here $\pi$ is a uniformizer, $i$ is unramified lifting, $a,b\in\Br(k)\{\ell\}$, and $\chi,\psi\in H^1(k,\Q_\ell/\Z_\ell)$.
This decomposition of unramifiedly split classes holds at every prime and respects residue-field restriction for finite extensions of ramification index one \cite[\S3, (3.3)--(3.4), Theorem 3.2]{AljadeffSonnWadsworth}.
Put $m=\operatorname{ord}\chi$ and $L=k(\chi)$, with $m=1,L=k$ when $\chi=0$.
Every finite $M/k$ lifts to a finite extension $F_M/F$ of ramification index one and residue field $M$: lift monic minimal polynomials successively along a tower of simple extensions.
The lifted ring is a complete discrete valuation ring, since its residue is a field and its maximal ideal is generated by $\pi$.
For every finite splitting field $M/L$ of $a_L$, both $a_M$ and $\chi_M$ vanish, since $M$ splits $a_L$ and contains $L$, so $\alpha_{F_M}=0$ by the Witt decomposition.
Since $\beta\in B_\alpha(F)$, also $\beta_{F_M}=0$, and uniqueness of the Witt decomposition over $F_M$ gives $\psi_M=0$.
By \cite[Proposition 5]{KatoExistence} with $n=q=h$ and $r=2$, we have $\sum_M N_{M/L}K_h^M(M)=K_h^M(L)$, where $M$ runs over these splitting fields.
The projection formula and duality on the $H^1(L,\Q_\ell/\Z_\ell)$ side therefore give $\psi_L=0$, hence $\psi=c\chi$ for some integer $c$ \cite[Proposition 5, Theorem 2]{KatoExistence}.
Replacing $\beta$ by $\beta-c\alpha$, we may assume $\beta=i(b)$.

We take any $u\in k^\times$ and any finite splitting field $E/k$ of $z=a+(\chi,u)$.
We adjoin $t^m=\pi/\tilde u$ to the lift $F_E$, where $\tilde u$ is a unit lift of $u$.
For $m>1$ this is an Eisenstein extension, so the residue remains $E$.
Over this field, writing $\pi=\tilde u\,t^m$, we have $(\chi,\pi)=(\chi,\tilde u)+m(\chi,t)=i((\chi,u))$ because $m\chi=0$; hence $\alpha=i(z_E)=0$ there, and so $i(b_E)=0$.
Unramified lifting is injective, so $b_E=0$.
The hypothesis and Lemma~\ref{lem:second-three-cyclic-family}, applied over $k$ with $\eta=a$ and $\theta=\chi$ of order $n=m$, give
\[
 b\in\bigcap_{u\in k^\times}\langle a+(\chi,u)\rangle
   =\langle ma\rangle,\qquad i(ma)=m\alpha .
\]
Thus the original $\beta$ is an integral multiple of $\alpha$.
\end{proof}

\begin{theorem}\label{thm:second-three-coker}
Let $F$ be a three-dimensional local field, let $D$ be a central division algebra over $F$, and let $\alpha$ be the Brauer class of $D$.
Then
\[
\begin{aligned}
 N_2(\alpha)=\Nrd_2K_2(D)
       &=\ker(\alpha\cup-:K_2(F)\to H^4(F,\Q/\Z(3))),\\
 C_2(D/F)&\xrightarrow[\ \inv_F(\alpha\cup-)\ ]{\ \sim\ }
       \tfrac{1}{\per(\alpha)}\Z/\Z .
\end{aligned}
\]
\end{theorem}

We prove Theorem~\ref{thm:second-three-coker} at the end of this subsection, after a sequence of lemmas.
For a three-dimensional local field $F$ and $\alpha\in\Br(F)$, we write $B_\alpha$ for $B_\alpha(F)$.
Lemma~\ref{lem:second-three-duality} reduces the theorem to the equality $B_\alpha=\langle\alpha\rangle$, and Lemma~\ref{lem:second-three-containment} in Section~\ref{sec:second-two-tame} places the image of the reduced norm in the kernel of the cup product.
For a wild class of period $p$, the remaining lemmas pass to an extension over which the relevant classes stay nonzero and $\alpha$ has index $p$ (Lemma~\ref{lem:second-three-wild-descent}), and they give the norm image there (Lemmas~\ref{lem:second-three-kahn-levine} and~\ref{lem:second-three-char-p-index-p}).
In the proof of the theorem, classes split by an unramified extension are treated first by Lemma~\ref{lem:second-three-witt-common-kernel}, and the equality $B_\alpha=\langle\alpha\rangle$ over $F$ for wild $p$-primary classes of higher period is proved by induction.

\begin{lemma}\label{lem:second-three-duality}
Let $F$ be a three-dimensional local field and $\alpha\in\Br(F)$.
Then $N_2(\alpha)$ is an open subgroup of finite index in $K_2(F)$, and Kato duality gives
\[
 (K_2(F)/N_2(\alpha))^\vee=B_\alpha .
\]
Moreover, $B_\alpha=\langle\alpha\rangle$ if and only if $N_2(\alpha)=\ker(\alpha\cup-)$, and in this case $\inv_F(\alpha\cup-)$ induces an isomorphism $K_2(F)/N_2(\alpha)\simeq\tfrac1{\per(\alpha)}\Z/\Z$.
\end{lemma}
\begin{proof}
The subgroup $N_2(\alpha)$ of $K_2(F)$ is open and of finite index \cite[Theorems 2--3]{KatoExistence}, so all characters of $K_2(F)/N_2(\alpha)$ are continuous.
Under Kato duality, a continuous character of $K_2(F)$ has the form $x\mapsto\inv_F(\beta\cup x)$ for a unique Brauer class $\beta$.
For each finite splitting field $E/F$, separable or not, the projection formula and the compatibility of the top invariant with corestriction show that this character vanishes on $N_{E/F}K_2(E)$ if and only if $\beta_E=0$ \cite[\S3.2, Lemma 1(2) and Proposition 1(2)(ii)]{KatoLocalII}.
Thus the dual of the finite quotient $K_2(F)/N_2(\alpha)$ is $B_\alpha$.

The class $\alpha$ generates $B_\alpha$ if and only if its character is faithful.
Since $\inv_F$ is an isomorphism, this holds if and only if $N_2(\alpha)=\ker(\alpha\cup-)$, and then the image of the character is the subgroup of order $\per(\alpha)$.
\end{proof}

\begin{lemma}\label{lem:second-three-kahn-levine}
Let $B$ be a central simple algebra over a field $k$ whose index $e$ is squarefree and prime to $\operatorname{char}k$.
Then
\[
 K_2(B)=\sum_{\substack{E/k\ \mathrm{finite\ separable}\\ B_E\ \mathrm{split}}}\theta_EK_2(E),
\]
and $\Nrd_2K_2(B)=N_2([B])$ is the kernel of
\[
 [B]\cup-\colon K_2(k)\longrightarrow H^4\bigl(k,(\Z/e)(3)\bigr).
\]
\end{lemma}
\begin{proof}
By \cite[Proposition 6.6.1(2)]{KahnLevine2010}, $K_2(B)$ is generated by the images of $\theta_E$ for finite splitting fields $E/k$.
If a finite splitting field $E/k$ is not separable, let $E_s$ be its maximal separable subextension; the index of $B_{E_s}$ divides both $e$, which is prime to $\operatorname{char}k$, and $[E:E_s]$, a power of $\operatorname{char}k$, so $E_s$ splits $B$ and $\theta_E=\theta_{E_s}N_{E/E_s}$.
This gives the displayed generation, and consequently $\Nrd_2K_2(B)=N_2([B])$.

The inclusion of $N_2([B])$ in the kernel of $[B]\cup-$ follows from the projection formula.
For the reverse inclusion, suppose first that $k$ is perfect.
Then \cite[Theorem 8.2.2(4)]{KahnLevine2010}, which uses the Beilinson--Lichtenbaum comparison in weights up to three \cite{Voevodsky}, identifies $N_2([B])$ with the kernel of $[B]\cup-$ after the exponential characteristic is inverted.
Both quotients of $K_2(k)$, by $N_2([B])$ and by the kernel of $[B]\cup-$, are killed by $e$: the first because $N_{E/k}\res_{E/k}=e$ for a separable maximal subfield $E$ of degree $e$ of the division algebra in the class of $B$, and the second because the coefficients are $\Z/e$.
Since $e$ is prime to $\operatorname{char}k$, inverting the exponential characteristic does not change the comparison, so the identification holds integrally.
If $\operatorname{char}k=q>0$ and $k$ is not perfect, let $z$ lie in the kernel of $[B]\cup-$.
Over the perfect closure, $z$ lies in the norm subgroup; we descend the resulting norm presentation to a finite purely inseparable stage and transfer it back.
This puts a $q$-power multiple of $z$ into $N_2([B])$, and since $q$ is prime to $e$, the element $z$ itself lies in $N_2([B])$.
\end{proof}

\begin{lemma}\label{lem:second-three-coefficients}
Let $F'$ be a three-dimensional local field and $p$ a prime.
The natural map
\[
 H^4(F',\Z/p(3))\longrightarrow H^4(F',\Q_p/\Z_p(3))
\]
is injective, and the top invariant identifies $H^4(F',\Q_p/\Z_p(3))$ with $\Q_p/\Z_p$.
If $p=\operatorname{char}F'$, the coefficients are logarithmic.
\end{lemma}
\begin{proof}
By the coefficient exact sequence, the kernel is the cokernel of multiplication by $p$ on $H^3(F',\Q_p/\Z_p(3))$.
The norm residue isomorphism in degree three identifies this group with the $p$-divisible group $K_3^M(F')\otimes\Q_p/\Z_p$: for $p\ne\operatorname{char}F'$ by \cite[Theorems 6.16 and 6.17]{Voevodsky}, and for $p=\operatorname{char}F'$ by \cite[Corollary (2.8)]{BlochKato1986}.
For $p=\operatorname{char}F'$, successive Kato residue maps followed by the trace on the finite last residue field identify $H^4(F',\Q_p/\Z_p(3))$ with $\Q_p/\Z_p$, and the top invariant does the same in general \cite[\S5.2]{Kurihara2000}.
\end{proof}

\begin{lemma}\label{lem:second-three-char-p-index-p}
Let $F'$ be a three-dimensional local field of characteristic $p$, let $B$ be a central simple algebra over $F'$ of index $p$, and put $\beta=[B]$.
Then
\[
 N_2(\beta)
 =\ker\bigl(\beta\cup-\colon K_2(F')
       \longrightarrow H^4(F',\Q_p/\Z_p(3))\bigr).
\]
\end{lemma}
\begin{proof}
The inclusion from left to right follows from the projection formula in Galois cohomology \cite[\S2 (C3)]{Otabe2023}, applied with finite coefficients and passed to the direct limit: for any finite splitting field $E/F'$ and any $w\in K_2(E)$,
\[
 \beta\cup N_{E/F'}w
 =\cor_{E/F'}(\beta_E\cup w)=0.
\]
For the reverse inclusion, take $z\in K_2(F')$ with $\beta\cup z=0$.
By Lemma~\ref{lem:second-three-coefficients}, $H^4(F',\Z/p(3))\simeq\Z/p$ is identified with the subgroup of $H^4(F',\Q_p/\Z_p(3))$ killed by $p$.
Let $\overline c=\beta\cup z\in H^4(F',\Z/p(3))$ be the product of $\beta\in\Br(F')[p]$ with the norm-residue symbol of $z$.
Its image in the direct-limit coefficient group is zero, so $\overline c=0$ by this identification.

We choose a complete Cohen discrete valuation ring $R$ of mixed characteristic $(0,p)$ with residue field $F'$, uniformizer $\pi$, and fraction field $\mathcal K$.
The Henselian Brauer-group isomorphism \cite[IV, Corollary 2.13]{MilneEtale}
\[
 \Br(R)\xrightarrow{\ \sim\ }\Br(F')
\]
and the fixed-degree lifting theorem for $\operatorname{PGL}_p$-torsors over Henselian local rings \cite[Expos\'e XXIV, Proposition 8.1]{SGA3} provide an Azumaya algebra $\widetilde B$ over $R$ of degree $p$ whose special fiber is isomorphic to the division algebra Brauer equivalent to $B$.
Its generic fiber $\widetilde B_{\mathcal K}$ also has index $p$: the index divides its degree, and the Brauer class cannot vanish over $\mathcal K$ because $\Br(R)\to\Br(\mathcal K)$ is injective \cite[IV, Corollary 2.6]{MilneEtale}.
By Matsumoto's theorem, we write $z$ as a finite sum of Milnor symbols; lifting their entries to units of $R$ gives
\[
 \widetilde z\in K_2^M(R),\qquad
 \operatorname{sp}(\widetilde z)=z.
\]

We use Kato's unramified lifting map with finite $p$-coefficients
\[
 \iota_{F'/\mathcal K}^q\colon
 H^q(F',\Z/p(q-1))
 \longrightarrow H^q(\mathcal K,\mu_p^{\otimes(q-1)})
\]
for $q=2,4$.
It is defined by lifting Artin--Schreier characters and the unit entries of symbols \cite[Theorem~3]{KatoGaloisCohomology1982}, \cite[\S5.1]{Kurihara2000}.
In degree two, we write $\beta$ as a sum of cyclic symbols and lift their Artin--Schreier characters to finite \'etale $R$-algebras and their units to $R^\times$.
The sum of the lifted cyclic classes extends over $R$ with special fiber $\beta$, so $\Br(R)\simeq\Br(F')$ gives $\iota_{F'/\mathcal K}^2(\beta)=[\widetilde B_{\mathcal K}]$.
In degree four, Kato's explicit product formula gives
\[
 \iota_{F'/\mathcal K}^4(\beta\cup z)
  =[\widetilde B_{\mathcal K}]\cup\widetilde z_{\mathcal K}
  \quad\text{in }H^4(\mathcal K,\mu_p^{\otimes3}).
\]
The left-hand side is zero because $\overline c=0$, so the finite-coefficient cup product over $\mathcal K$ vanishes.

Since $\mathcal K$ has characteristic zero and $\widetilde B_{\mathcal K}$ has index $p$, Lemma~\ref{lem:second-three-kahn-levine} gives finite separable splitting fields $L_j/\mathcal K$ of $\widetilde B_{\mathcal K}$ and elements $w_j\in K_2(L_j)$ such that
\[
 \widetilde z_{\mathcal K}=\sum_jN_{L_j/\mathcal K}(w_j)
    \quad\text{in }K_2(\mathcal K).
\]
The integral closure $R_j$ of $R$ in $L_j$ is a complete discrete valuation ring; let $F'_j$ be its residue field.
Since $\Br(R_j)\to\Br(L_j)$ is injective and $\widetilde B_{L_j}$ splits, each finite residue extension $F'_j/F'$ splits $B$.
The residue extensions need not be separable, which is harmless because the definition of $N_2(\beta)$ includes all finite splitting fields.

Write $\partial_{\mathcal K}\colon K_3(\mathcal K)\to K_2(F')$ and $\partial_{L_j}\colon K_3(L_j)\to K_2(F'_j)$ for the boundary maps in Quillen's localization sequences, and identify $K_2$ of the residue fields with Milnor $K_2$ by Matsumoto's theorem; under this identification the Quillen and Milnor norms agree.
Since $\widetilde z$ comes from $K_2(R)$, the $K(R)$-linearity of the boundary and $\partial_{\mathcal K}\{\pi\}=1$ give the first equality below.
The projection formula in Quillen $K$-theory, which expresses the compatibility of tensor products with restriction of scalars, and the compatibility of norms with boundaries then give
\[
 \begin{aligned}
 z
 &=\partial_{\mathcal K}\bigl(\{\pi\}\cdot\widetilde z_{\mathcal K}\bigr)\\
 &=\sum_j\partial_{\mathcal K}N_{L_j/\mathcal K}
             \bigl(\res_{L_j/\mathcal K}\{\pi\}\cdot w_j\bigr)\\
 &=\sum_jN_{F'_j/F'}\,\partial_{L_j}
             \bigl(\res_{L_j/\mathcal K}\{\pi\}\cdot w_j\bigr).
 \end{aligned}
\]
The compatibility of norms with boundaries in the last equality follows from Quillen's localization sequence and d\'evissage \cite[\S5, Theorems 4 and 5]{QuillenHigher}: restriction of scalars along the finite free extension $R_j/R$ preserves modules of finite length, and on closed support it is restriction of scalars from $F'_j$ to $F'$.
Hence the boundary commutes with norms even when $F'_j/F'$ is inseparable, and no ramification factor appears.
Since each $F'_j$ splits $B$, every term of the final sum lies in $N_2(\beta)$, and hence $z\in N_2(\beta)$.
\end{proof}

The remaining lemmas treat wild classes of period $p$ by index descent.
In both, $F$ is a three-dimensional local field with first residue characteristic $p$ and first residue field $k$, $\alpha\in\Br(F)[p]$ is wild, and $S\subset\Br(F)[p]\setminus\{0\}$ is a finite set containing $\alpha$ with at most one purely unramified class.
In mixed characteristic we first adjoin $\mu_p$ and, if necessary, pass to a totally ramified extension of degree prime to $p$; these restrictions are injective on $p$-torsion and preserve wildness, and they create no new purely unramified classes, since the character component of a Witt class becomes $e\cdot\res\chi$ with $e$ prime to $p$.
We choose the degree of the totally ramified extension to achieve the depth inequalities used below.
After these preparatory extensions we again write $F$, $\alpha$ and $S$ for the field and the classes, and only then choose $k$, the uniformizer $\pi$ and the residue $p$-basis.

\begin{lemma}\label{lem:second-three-lambda}
In this situation, there are a residue $p$-basis $(a,b)$ of $k$ and a unit $C$ of $F$ with residue $a$ such that, for $F_i=F(C^{1/p^i})$ and the completion $L^\wedge$ of $L=\bigcup_iF_i$, every member of $S$ is nonzero over $L^\wedge$ and $\alpha_{L^\wedge}$ is wild.
The residue field $k_\infty$ of $L^\wedge$ satisfies $[k_\infty:k_\infty^p]=p$, with $p$-basis $b$.
\end{lemma}
\begin{proof}
For $\gamma\in\Br(F)[p]$ we use the leading term of $\gamma$ for Kato's filtration in equal characteristic (indexed by the pole order $i$) and for the Bloch--Kato filtration in mixed characteristic (indexed by the depth $j$ of $U^j$).
In equal characteristic for pole order $i>0$, and in mixed characteristic for depth $0<j<c$, its graded pieces are $\Omega^1_k$ if $p$ does not divide the index and $\Omega^1_k/Z^1_k\oplus k/k^p$ if $p$ divides it; in mixed characteristic, the piece in degree zero is $\nu^2(k)\oplus\nu^1(k)$.
The terminal layer, pole order zero in equal characteristic and depth $c$ (defined below) in mixed characteristic, is the Witt part.
The components below are those of this leading term, and a class is wild exactly when its leading term lies outside the terminal layer.
For a $p$-basis $(a,b)$ of $k$, $\partial_a,\partial_b$ are the derivations dual to $da,db$.

We perfect one direction of a residue $p$-basis $(a,b)$, keeping every member of $S$ nonzero and keeping $\alpha$ wild.
The forbidden directions for $da$ are the nonzero one-form components of $S$, the differentials $dw$ of nonzero components $w\in k/k^p$, and the logarithmic-differential directions of nonzero degree-zero $\nu^1(k)$ components; there are finitely many of them.
A nonzero character component in the Witt decomposition remains nonzero under purely inseparable extension.
If $S$ contains no purely unramified class, we choose $a$ directly so that $da$ avoids the forbidden directions.
Otherwise, let $\gamma\in\Br(k)[p]$ be the purely unramified class in $S$.
Two-dimensional local duality gives $a_0$ with $\gamma\cup a_0\ne0$.
Then $a_0\notin k^p$, so $da_0$ extends to a basis of $\Omega^1_k$.
We choose $db_0$ independent of $da_0$, and vary $a=a_0+t b_0$ with sufficiently deep $t\in k^p$.
The cup-product character kills some principal-unit group, so $\gamma\cup a\ne0$ is preserved.
Meanwhile, $da=da_0+t\,db_0$ takes infinitely many directions and can avoid the finitely many forbidden ones.
For $k_i=k(a^{1/p^i})$, the projection formula and $N_{k_i/k}(a^{1/p^i})=a$ give $\gamma_{k_i}\ne0$ at every stage, and hence over the direct limit $k_\infty$.
We choose $b$ to complete $a$ to a $p$-basis.

We lift $a$ to a unit $C$, and put $F_i=F(C^{1/p^i})$.
The polynomial degree gives $[F_i:F]\le p^i=[k_i:k]$, and the residue field contains $k_i$, so $F_i/F$ has ramification index one and residue field $k_i$.
The residue $k_\infty$ of the completion $F_\infty^\wedge=L^\wedge$ of $\bigcup_i F_i$ has $p$-basis $b$.
In particular, $\ker(\Omega^1_k\to\Omega^1_{k_\infty})=k\,da$.
We combine this with the freedom to vary $C$ modulo $\pi^2$.

Suppose first that $F$ has equal characteristic, and write $F=k((\pi))$.
Kato's leading term at pole order $i>0$ is measured by $\Omega^1_k$ if $p\nmid i$, and by $\Omega^1_k/Z^1_k\oplus k/k^p$ if $p\mid i$, where $Z^1_k=\ker d$.
In the first case, a nonzero form remains nonzero if the direction of $da$ avoids that form; in the second case, a nonzero second component $w\bmod k^p$ remains nonzero if the direction of $da$ avoids $dw$.
If the second component vanishes, take a representative $\omega=f\,da+g\,db$ of a nonzero class in $\Omega^1_k/Z^1_k$; then $d\omega\ne0$.
We set $C=a+\lambda\pi$ for $\lambda\in k^p$ and choose a new coefficient field of $F_\infty^\wedge$ for its residue field $k_\infty$ containing the prescribed lift of $b$; let $a^\sharp$ be the lift of $a$ in it.
Then $C$ and $a^\sharp$ have the same residue and both have all $p$-power roots in $F_\infty^\wedge$, so if they differed, the positive valuation of $C/a^\sharp-1$ would be divisible by every $p^n$; hence $C=a^\sharp$.
To first order, the difference between the old and new lifts of any $u\in k$ is $-\lambda\pi\partial_a u$, by the binomial expansion of $u=\sum u_{ij}^p a^i b^j$.
Since $dC=0$ in $F_\infty^\wedge$, we have
\[
 \omega\equiv g\,db-\lambda\pi
       \bigl((\partial_a g)\,db+f\,d\log\pi\bigr)
       \pmod{\pi^2}.
\]
In this case, where $p\mid i$, the second component vanishes and the representative satisfies $d\omega\ne0$, the closed form lowered by Cartier to pole order $i/p$ is the restricted form $g\,db$; it is closed over the coefficient field, which has a one-element $p$-basis.
At pole order $i-1$, eliminating the $d\log\pi$ term with the exact differential $d(\pi^{-(i-1)}\lambda f)$ makes the coefficient of $\lambda$ equal to $(\partial_b f-\partial_a g)\,db$, which is nonzero since $d\omega\ne0$.
The coefficient of the image of the restricted closed form $g\,db$ under the Cartier operator is given by its residue in the new coefficient field, which does not involve $\lambda$.
This also applies when $i/p=i-1$, which happens only for $i=p=2$.

Suppose next that $F$ has mixed characteristic, so that the preliminary extensions above have been made.
Put $e=v_F(p)$ and $c=pe/(p-1)$.
The terminal term $U^c$ of the Bloch--Kato unit filtration modulo $p$ corresponds to the Witt part: with $\tau=\zeta_p-1$, the residue of $\bigl((1+\tau y)^p-1\bigr)/\tau^p$ is $y^p-y$.
Thus the Kummer character of $1+(\zeta_p-1)^pu$ is the unramified lift of the Artin--Schreier character of $\bar u$, and taking a unit and a uniformizer in the second slot generates the two Witt components \cite[(5.1)--(5.3)]{BlochKato1986}.
For $0<j<c$, the graded pieces are $\Omega^1_k$ if $p\nmid j$ and $\Omega^1_k/Z^1_k\oplus k/k^p$ if $p\mid j$.
The degree-zero piece is $\nu^2(k)\oplus\nu^1(k)$, where $\nu^r(k)=K_r^M(k)/p$ is identified by logarithmic differentials \cite[(5.1)--(5.3)]{BlochKato1986}.
Scaling $e,c,j$ by the degree of the preparatory totally ramified extension, we have arranged $j+1<c$ for each relevant positive depth, and $e+j/p-j\ge2$ when $p\mid j$; we have also arranged $e=v_F(p)\ge2$, which is used below even when $S$ has no class of positive depth.
For the change $\pi'^d=\pi$ with $p\nmid d$, we have $d\log\pi=d\cdot d\log\pi'$, so the nonzero graded pieces are preserved while these depth inequalities are achieved.
The choice of $da$ avoiding the forbidden directions preserves nonzero one-forms (for $p\nmid j$), components in $k/k^p$, and nonzero degree-zero $\nu^1(k)$ components.

For the remaining cases, take integral lifts $\widetilde a,\widetilde b,\widetilde\lambda$ of $a,b,\lambda$ and put $C=\widetilde a+\widetilde\lambda\pi$.
Fix a lift $\widetilde u$ of $u=\sum x_{rs}^p a^r b^s$, and set $\widetilde u_0=\sum \widetilde x_{rs}^pC^r\widetilde b^s$.
Then
\[
 \overline{(\widetilde u-\widetilde u_0)/\pi}=z_{\widetilde u}-\lambda\,\partial_a u,
\]
where $z_{\widetilde u}$ is the value of the left-hand side at $C=\widetilde a$.
For positive depth $j$ divisible by $p$ with zero second component, lift $\omega=f\,da+g\,db$ by $\{1+\pi^j\widetilde f \widetilde a,\widetilde a\}+ \{1+\pi^j\widetilde g \widetilde b,\widetilde b\}$.
In the first term, $C$ in $\widetilde a=C(1-\widetilde\lambda\pi/C)$ becomes a $p$-th power.
We apply the first-order Steinberg relation $\{1+x,1+y\}\equiv\{1+xy,-x^{-1}\}\bmod U^{j+2}$ \cite[(4.1.1)]{BlochKato1986} to the first symbol with $x=\pi^j\widetilde f\widetilde a$ and $y=-\widetilde\lambda\pi/C$.
Then $\overline{xy/\pi^{j+1}}=-\lambda f$, and since $p\mid j$ and $da=0$ after restriction, the term of degree $j+1$ in $-\lambda f\,d\log(-x^{-1})$ is $\lambda\,df$.
The $\widetilde u_0$-part of the coefficient of the second term can be written modulo $\pi^2$ as $\sum_{r<p}h_r^p \widetilde b^r$.
Since $e+j/p-j\ge2$, we have $1+\pi^jh_0^p\equiv(1+\pi^{j/p}h_0)^p\pmod{U^{j+2}}$, so the $r=0$ part is a $p$-th power modulo $U^{j+2}$.
For $r>0$, multiplying the symbol by $r$ and using Steinberg moves it to a symbol with second slot $-(\pi^j h_r^p)^{-1}$, and the same vanishing follows.
The new second slot of the second symbol is a $p$-th power up to sign, and the sign for $p=2$ vanishes in this degree since $-1\in U_F^e$ and $e\ge2$.
Thus $-\lambda(\partial_a g)\,db$ is the $\lambda$-dependent part of the second term; a term independent of $\lambda$ may also occur, which is why at most one value of $\lambda$ per class is excluded below.
Hence the coefficient of $\lambda$ is the nonzero form $(\partial_b f-\partial_a g)\,db$.

At degree zero, if the $\nu^1(k)$ component vanishes, lift the nonzero $\omega\in\nu^2(k)$ by a sum of unit symbols.
The ring $\mathcal O_{F_\infty^\wedge}/(\pi^2)$ has characteristic $p$.
We consider the section sending an element of $k_\infty^p$ to the $p$-th power of a lift of its root; modulo $\pi^2$ in characteristic $p$, this $p$-th power does not depend on the lift and is compatible with $p$-th powers of $b$, so the section is a ring homomorphism.
Using the power basis in $b$, we extend this section to a section of $k_\infty$.
Symbols along this section vanish modulo $(p,U^2)$, since $K_2^M(k_\infty)/p=0$.
Taking the difference in the preceding first-order expansion, we find that the coefficient of $\lambda$ for an original symbol $\{\widetilde u,\widetilde v\}$ is
\[
 -(\partial_a u/u)\,d\log v+(\partial_a v/v)\,d\log u.
\]
For the sum it is therefore $-\iota_{\partial_a}\omega$; a nonzero $da\wedge db$ component maps to a nonzero $db$ component.

In every case, contributions already lying in the next layer come from residue restriction and are independent of $\lambda$.
Avoiding at most one $\lambda$ for each class keeps all members of $S$ nonzero and $\alpha$ wild.
\end{proof}

The following density statement is used in the proof of the next lemma.

\begin{lemma}\label{lem:paper-henselian-density}
Let $H$ be a Henselian field for a rank-$1$ valuation, and let $\widehat H$ be its completion.
On any smooth finite-type $H$-scheme, the $H$-points are dense in the $\widehat H$-points.
\end{lemma}
\begin{proof}
Locally, a smooth scheme is standard \'etale over $\mathbb A^n_H$: it is given by $g(a,y)=0$ with $\partial_yg(a,y)\ne0$, where $g\in H[a_1,\dots,a_n][y]$ is monic in $y$.
Let $(a,b)$ be an $\widehat H$-point.
After translating and rescaling, we may assume that the coefficients of $g(a,-)$ and the point lie in the valuation ring.
Since $H$ is dense in $\widehat H$, we may choose $a'\in H^n$ and $b'\in H$ so close to $a$ and $b$ that $|g(a',b')|<|\partial_yg(a',b')|^2$.
Newton's form of Hensel's lemma in the Henselian rank-one field $H$ then gives a root $b''\in H$ of $g(a',-)$ with $|b''-b'|\le|g(a',b')|/|\partial_yg(a',b')|$, which can be made arbitrarily small, and $(a',b'')$ is an $H$-point close to $(a,b)$.
Shrinking the neighborhood keeps the open condition $\partial_yg\ne0$.
\end{proof}

\begin{lemma}\label{lem:second-three-wild-descent}
In this situation, there is a finite extension $F'/F$, namely a stage $F_i$ of the tower in Lemma~\ref{lem:second-three-lambda}, such that every member of $S$ is nonzero over $F'$, $\alpha_{F'}$ is wild, and $\ind(\alpha_{F'})=p$.
\end{lemma}
\begin{proof}
Let $(a,b)$, $C$, $F_i$ and $L$ be as in Lemma~\ref{lem:second-three-lambda}.
The residue field $k_\infty$ of $L^\wedge$ satisfies $[k_\infty:k_\infty^p]=p$, so the wild class $\alpha_{L^\wedge}$ of period $p$ has index $p$ by \cite[\S4, Lemma 5, p.~337]{KatoLocalI}, whose hypotheses (a $p$-torsion class not split by the maximal unramified extension, and residue $p$-rank one) hold here.
Let $A$ be a central simple algebra over $F$ with class $\alpha$, and let $X=\SB_p(A_L)$ be the generalized Severi--Brauer variety parametrizing right ideals of reduced dimension $p$.
Since $X$ is smooth and $L$ is Henselian and dense in $L^\wedge$, Lemma~\ref{lem:paper-henselian-density} and $X(L^\wedge)\ne\varnothing$ give $X(L)\ne\varnothing$.
Finite presentation descends this point to a finite stage $F_i$, where $\ind(\alpha_{F_i})\mid p$.
All members of $S$ are nonzero over $L^\wedge$, hence over $F_i$; in particular, $\ind(\alpha_{F_i})=p$.
If $\alpha_{F_i}$ belonged to the Witt part, so would its restriction $\alpha_{L^\wedge}$, because a finite unramified extension splitting a Witt class lifts along the base change of ramification index one; this contradicts Lemma~\ref{lem:second-three-lambda}.
Hence $\alpha_{F_i}$ is wild.
\end{proof}

\begin{proof}[Proof of Theorem~\ref{thm:second-three-coker}]
By Lemma~\ref{lem:second-three-duality}, $(K_2(F)/N_2(\alpha))^\vee=B_\alpha$.
The splitting transfers give $N_2(\alpha)\subset\Nrd_2K_2(D)$, and Lemma~\ref{lem:second-three-containment} gives $\Nrd_2K_2(D)\subset\ker(\alpha\cup-)$.
By definition, $B_\alpha$ contains $\langle\alpha\rangle$, and it remains to prove equality.
Since $B_\alpha=\bigoplus_\ell B_{\alpha_\ell}$ and $\langle\alpha\rangle=\bigoplus_\ell\langle\alpha_\ell\rangle$, it suffices to do so for $\ell$-primary $\alpha$.
A class in the Witt part is split by a finite unramified extension, so Lemma~\ref{lem:second-three-witt-common-kernel} applies with $h=2$, its hypothesis being the two-dimensional case recalled at the beginning of this subsection; this covers all $\ell$-primary classes with $\ell$ different from the first residue characteristic, since they lie in the Witt part.

Suppose that $\alpha$ is wild of period $p$ and that $\beta\in B_\alpha[p]\setminus\langle\alpha\rangle$.
Put $S=\{\alpha\}\cup\{\beta-j\alpha:j\in\F_p\}$.
Every member of $S$ is nonzero, and $S$ contains at most one purely unramified class: the difference of two would be a nonzero multiple of the wild class $\alpha$.
Lemma~\ref{lem:second-three-wild-descent} gives a finite extension $F'=F_i$ of the original field $F$, through the preparatory extensions, over which every member of $S$ is nonzero, $\alpha_{F_i}$ is wild, and $\ind(\alpha_{F_i})=p$.
We claim that $N_2(\alpha_{F_i})=\ker(\alpha_{F_i}\cup-)$.
In characteristic zero, apply Lemma~\ref{lem:second-three-kahn-levine} with $k=F_i$ and $e=p$, and use Lemma~\ref{lem:second-three-coefficients} to pass from $\Z/p(3)$ to $\Q/\Z(3)$ coefficients.
In characteristic $p$, apply Lemma~\ref{lem:second-three-char-p-index-p} with $F'=F_i$ and $\beta=\alpha_{F_i}$.
By the equivalence in Lemma~\ref{lem:second-three-duality} over $F_i$, the claim gives $B_{\alpha_{F_i}}=\langle\alpha_{F_i}\rangle$.
Since $\beta_{F_i}\in B_{\alpha_{F_i}}$, there is $j\in\F_p$ such that the element $\beta-j\alpha$ of $S$ vanishes over $F_i$, contradicting the nonvanishing of the members of $S$.

Next put $\per(\alpha)=p^r$ and $\theta=p^{r-1}\alpha$.
For any finite extension $E/F$ splitting $\theta$, let $E_s$ be its maximal separable subextension.
In the Galois closure $E^{\mathrm{gal}}_s$ of $E_s/F$, choose a Sylow $p$-subgroup $P$ of $\Gal(E^{\mathrm{gal}}_s/F)$ containing a Sylow $p$-subgroup of $\Gal(E^{\mathrm{gal}}_s/E_s)$, and put $F''=(E^{\mathrm{gal}}_s)^P$.
Both $[F'':F]$ and $[EF'':E]$ are prime to $p$; for the latter, $[E_sF'':E_s]=[H:H\cap P]$ with $H=\Gal(E^{\mathrm{gal}}_s/E_s)$, and $E/E_s$ is purely inseparable.
A subgroup chain of the $p$-group gives a tower of degree-$p$ extensions for $E_sF''/F''$; the purely inseparable extension $EF''/E_sF''$ also has a degree-$p$ tower.
Let $M'/M$ be the first step where $\theta$ vanishes.
We have $\per(\alpha_M)=p^r$ and $\per(\alpha_{M'})=p^{r-1}$.
Restriction in degree $p$ lowers order by at most a factor of $p$: for a separable step this follows from corestriction, and for a purely inseparable step $M'/M$ of degree $p$, the composite $M\hookrightarrow M'\xrightarrow{x\mapsto x^p}M$ is the Frobenius, which acts on $\Br(M)\{p\}$ as multiplication by $p$.
For $r\ge2$, we assume by induction that $B_{\alpha'}=\langle\alpha'\rangle$ for every wild class $\alpha'$ with $p^{r-1}\alpha'=0$ over a three-dimensional local field.
By Lemma~\ref{lem:second-three-witt-common-kernel}, the same equality holds for every $p$-primary class $\alpha'$ in the Witt part over such a field.
Still for $r\ge2$, we take $\beta\in B_\alpha[p]$.
Every finite extension of $M'$ splitting $\alpha_{M'}$ also splits $\alpha$ over $F$; hence restriction sends $B_\alpha$ into $B_{\alpha_{M'}}$, and $\beta_{M'}\in B_{\alpha_{M'}}[p]$.
Induction gives $\beta_{M'}=j p^{r-2}\alpha_{M'}$ for some integer $j$.
If $j\not\equiv0\pmod p$, then $\beta_M-jp^{r-2}\alpha_M$ has order $p^2$ and cannot vanish in a degree-$p$ extension.
Thus $\beta_{M'}=0$, hence $\beta_{EF''}=0$.
As $[EF'':E]$ is prime to $p$, also $\beta_E=0$, so $\beta\in B_\theta[p]$.
Moreover, $B_\theta[p]=\langle\theta\rangle$, by Lemma~\ref{lem:second-three-witt-common-kernel} if $\theta$ lies in the Witt part, and by the period-$p$ case if $\theta$ is wild.
For $r=1$, we have already proved $B_\alpha[p]=\langle\alpha\rangle$.
In both cases, $B_\alpha$ is finite and $B_\alpha[p]$ has $\F_p$-dimension one, so $B_\alpha$ is cyclic.
If a generator $\delta$ had order $p^{r+q}$, $q>0$, then $\alpha$ would be a unit multiple of $p^q\delta$.
At the first vanishing step $M'/M$ for $\theta$, $\alpha_M$ has order $p^r$, so $\delta_M$ has order $p^{r+q}$.
Consequently, $\delta_{M'}$ has order at least $p^{r+q-1}$.
By the same restriction property as above, every finite extension of $M'$ splitting $\alpha_{M'}$ splits $\alpha$, so $\delta_{M'}\in B_{\alpha_{M'}}$.
But for $r\ge2$, by induction, $B_{\alpha_{M'}}=\langle\alpha_{M'}\rangle$ is killed by $p^{r-1}$.
For $r=1$, one has $\alpha_{M'}=0$ and hence $B_{\alpha_{M'}}=0$.
Both cases are contradictions.
Therefore $B_\alpha=\langle\alpha\rangle$, so Lemma~\ref{lem:second-three-duality} gives $N_2(\alpha)=\ker(\alpha\cup-)$.
With the inclusions $N_2(\alpha)\subseteq\Nrd_2K_2(D)\subseteq\ker(\alpha\cup-)$ shown above, this gives $N_2(\alpha)=\Nrd_2K_2(D)=\ker(\alpha\cup-)$, and the same lemma identifies $C_2(D/F)=K_2(F)/N_2(\alpha)$ with $\tfrac1{\per(\alpha)}\Z/\Z$ via $\inv_F(\alpha\cup-)$, which proves the theorem.
\end{proof}
